\begin{tikzpicture}[font=\figurefont\small,
 agent/.style={circle,draw=BookRule,fill=FigureInput!12,minimum size=10mm,align=center},
 box/.style={draw=BookRule,fill=FigureModel!12,minimum width=24mm,minimum height=9mm,align=center},
 flow/.style={-{Latex},draw=BookInk,line width=.8pt}]
 \node[agent] (a) at (1,0) {1};
 \node[agent] (b) at (-.4,-1.5) {2};
 \node[agent] (c) at (2.4,-1.5) {3};
 \draw[flow] (b)--node[left]{邻域边}(a);
 \draw[flow] (c)--(a);
 \node[above=2mm of a,text=BookTeal] {角色：通行者};
 \node[agent] (aa) at (7,0) {3};
 \node[agent] (bb) at (5.6,-1.5) {1};
 \node[agent] (cc) at (8.4,-1.5) {2};
 \draw[flow] (bb)--(aa);
 \draw[flow] (cc)--(aa);
 \node[above=2mm of aa,text=BookTeal] {角色：通行者};
 \draw[flow,dashed] (2.7,.1)--node[above]{仅重编号}(5.3,.1);
 \node[box] (agg) at (1,-3.4) {集合聚合\\$\operatorname{Agg}\{\bm m_{j\to1}\}$};
 \node[box] (aggr) at (7,-3.4) {同一聚合器\\$\operatorname{Agg}\{\bm m_{j\to3}\}$};
 \draw[flow] (b.south)--(agg.north west);
 \draw[flow] (c.south)--(agg.north east);
 \draw[flow] (bb.south)--(aggr.north west);
 \draw[flow] (cc.south)--(aggr.north east);
 \node[box,fill=FigureOutput!12] (out) at (1,-5) {$\pi_1(h_1,e_1,\bm z_1)$};
 \node[box,fill=FigureOutput!12] (outr) at (7,-5) {$\pi_3(h_3,e_3,\bm z_3)$};
 \draw[flow] (agg)--(out);
 \draw[flow] (aggr)--(outr);
 \draw[flow,dashed] (out)--node[above]{同一物理参与者的输出}(outr);
\end{tikzpicture}
